\section*{Bab \ref{ch:b1:rifts-ridges-margins} --- Celah, punggungan dan tepi pasif}

\begin{solution}{exo:b1:rifts-ridges-margins:1}
(a) Dike itu akan menutupi $\qty{60000}{m}\times\qty{8}{m} = \qty{4.8e5}{m^2}$ dalam denah,
sehingga tinggi rata-ratanya adalah $\qty{2.5e9}{m^3}/\qty{4.8e5}{m^2} =
\qty{5200}{m}$, sekitar \qty{5}{km}: dike itu memotong sebagian besar \omterm{def:b1:layered-earth:layers}{kerak}.
(b) $\qty{8}{m}/\qty{0.018}{m.yr^{-1}} = 444$ tahun, sekitar empat abad.
\end{solution}

\begin{solution}{exo:b1:rifts-ridges-margins:2}
(a) Pergeseran horizontal adalah lemparan dibagi tangen \omterm{def:b1:geological-maps:attitude}{kemiringannya}:
$3.0/\tan\ang{60} = \qty{1.73}{km}$ dan $0.8/\tan\ang{70} = \qty{0.29}{km}$;
perpanjangan penampang adalah \qty{2.02}{km}. (b) Sebelum pensesaran penampang itu
selebar $25 - 2.02 = \qty{22.98}{km}$, sekitar \qty{23}{km}: ia telah memanjang sekitar
\qty{9}{\percent}.
\end{solution}

\begin{solution}{exo:b1:rifts-ridges-margins:3}
Sangat lambat (di bawah \qty{20}{mm.yr^{-1}}): Punggungan Gakkel (12), Laut Merah tengah
(15) dan Punggungan Reykjanes (19), yang terakhir dekat dengan batas. Lambat (20 hingga
\qty{55}{mm.yr^{-1}}): Punggungan Carlsberg bagian barat (23) dan Atlantik Selatan (33).
Cepat (di atas \qty{75}{mm.yr^{-1}}): Punggungan Pasifik Timur (143). Punggungan yang
sangat lambat dan yang lambat seharusnya memiliki lembah aksial di antara sesar normal;
Punggungan Pasifik Timur memiliki tinggian aksial.
\end{solution}

\begin{solution}{exo:b1:rifts-ridges-margins:4}
(a) $\qty{3.0}{km^2}\times\qty{7.1}{km} = \qty{21.3}{km^3}$, sekitar
\qty{21}{km^3} setahun. (b) $\qty{3.028e8}{km^2}/\qty{3.0}{km^2.yr^{-1}} =
\qty{1.01e8}{yr}$, \qty{101}{Ma}. Dasar tertua, sekitar \qty{190}{Ma}, lebih tua karena
dasar samudra tidak dihancurkan menurut urutan umurnya: subduksi memakan dasar dari setiap
umur, sebagian dasar tua bertahan di tempat tidak ada palung yang mencapainya, dan laju
produksi telah bervariasi.
\end{solution}

\begin{solution}{exo:b1:rifts-ridges-margins:5}
(a) $2500 + 350\times2 = \qty{3200}{m}$, $2500 + 350\times6 = \qty{4600}{m}$,
$2500 + 350\times8 = \qty{5300}{m}$. (b) $t = \bigl((d - 2500)/350\bigr)^2$:
$(2100/350)^2 = \qty{36}{Ma}$ dan $(2800/350)^2 = \qty{64}{Ma}$. (c) Dasar tenggelam
\qty{1400}{m} antara 4 dan \qty{36}{Ma}, tetapi hanya \qty{700}{m} antara 36 dan
\qty{64}{Ma}, dalam kurun waktu yang kurang lebih sama.
\end{solution}

\begin{solution}{exo:b1:rifts-ridges-margins:6}
Koefisien muai efektif adalah
$\num{3.2e-5}\times3300/2300 = \qty{4.591e-5}{K^{-1}}$, dan
$\sqrt{\kappa\times\qty{1}{Ma}} = \sqrt{\num{e-6}\times\num{3.156e13}} =
\qty{5618}{m}$. Kemudian $C = (2/\sqrt\pi)\times\num{4.591e-5}\times1350\times5618 =
1.1284\times\num{4.591e-5}\times1350\times5618 = \qty{393}{m.Ma^{-1/2}}$. Untuk
\omterm{def:b1:layered-earth:layers}{mantel} yang \qty{250}{K} lebih dingin, $C$ berskala dengan $T_m - T_0$:
$393\times1100/1350 = \qty{320}{m.Ma^{-1/2}}$. Dasar samudra yang lebih dingin
tenggelam lebih lambat dan mengikuti konstanta yang lebih kecil. \num{350} yang dicocokkan
pada bab ini bukan kasus seperti itu: ia terletak dalam ketidakpastian parameternya.
\end{solution}

\begin{solution}{exo:b1:rifts-ridges-margins:7}
(a) $\qty{e-6}{cal.cm^{-2}.s^{-1}} = \num{4.184e-6}\times\num{e4} =
\qty{4.184e-2}{W.m^{-2}} = \qty{41.84}{mW.m^{-2}}$, sehingga konstantanya adalah
$11.3\times41.84 = \qty{473}{mW.m^{-2}.Ma^{1/2}}$, \qty{3.5}{\percent} di bawah
490 di \cref{ch:b1:heat-budget}: kedua kecocokan itu sepakat. (b) $q = 473$,
$473/\sqrt{10} = 150$ dan $473/\sqrt{150} = \qty{38.6}{mW.m^{-2}}$. Hukum akar kuadrat
terus turun, sedangkan model lempeng mendatar pada suatu konstanta: pada
\qty{150}{Ma} nilai akar kuadrat terlalu rendah, sebagaimana hukum kedalaman yang
terlalu dalam.
\end{solution}

\begin{solution}{exo:b1:rifts-ridges-margins:8}
Ambil segmen utara pada $x = 0$ dan segmen selatan pada $x = -40$
(km), dengan transform pada $y = 0$ di antara keduanya. (a) Tepat di utara transform,
dasar terletak di barat segmen utara dan bergerak ke barat; tepat di selatannya, dasar
terletak di timur segmen selatan dan bergerak ke timur. Dilihat dari kedua sisi, sisi
yang jauh bergerak ke kiri: geser-kirinya (sinistral). (b) Punggungan tampak tergeser
secara dekstral, bagian selatannya ke barat: arah yang berlawanan, seperti yang
diprediksi \cref{prop:b1:rifts-ridges-margins:transform}. (c) Gempa bumi hanya terjadi
di sepanjang \qty{40}{km} di antara kedua ujung punggungan itu, tempat kedua sisi
termasuk lempeng yang berbeda; di luarnya, sepanjang \omterm{def:b1:rifts-ridges-margins:transform}{zona patahan}, kedua sisi termasuk
lempeng yang sama dan bergerak bersama.
\end{solution}

\begin{solution}{exo:b1:rifts-ridges-margins:9}
Setengah laju adalah \qty{10.6}{mm.yr^{-1}}, \qty{10.6}{km.Ma^{-1}}. (a) Lompatan umur
adalah $350/10.6 = \qty{33.0}{Ma}$. (b) Pada ujung transform, dasar
berumur \qty{33.0}{Ma} ($2500 + 350\sqrt{33.0} = \qty{4511}{m}$) berhadapan dengan dasar
berumur 0 (\qty{2500}{m}): lompatan sekitar \qty{2010}{m}. \qty{100}{km} melewatinya,
umurnya $450/10.6 = 42.5$ dan $100/10.6 = \qty{9.4}{Ma}$, kedalamannya
4781 dan \qty{3575}{m}: lompatan sekitar \qty{1210}{m}. \qty{1000}{km} melewatinya,
umurnya 127.4 dan \qty{94.3}{Ma}, kedalamannya 6450 dan \qty{5899}{m}:
\qty{550}{m}. (Umur ini melampaui \qty{80}{Ma}, tempat hukum akar kuadrat menjadi
terlalu dalam; hukum lempeng akan memberi lompatan yang lebih kecil.)
\end{solution}

\begin{solution}{exo:b1:rifts-ridges-margins:10}
(a) Dengan luas yang sama untuk setiap umur, kedalaman rata-rata adalah rata-rata $d(t)$ atas
$[0, A]$:
\[
  \frac1A\int_0^A\bigl(d_r + C\sqrt t\bigr)\mathrm dt = d_r + \frac CA\cdot\frac23A^{3/2}
  = d_r + \frac23C\sqrt A.
\]
(b) Untuk $A = \qty{180}{Ma}$: $2500 + \frac23\times350\times13.416 =
\qty{5630}{m}$; untuk \qty{120}{Ma}: $2500 + \frac23\times350\times10.954 =
\qty{5056}{m}$, \qty{574}{m} lebih dangkal. (c) Misalkan permukaan laut naik sebesar $r$ dan
dasar yang terbebani tenggelam sebesar $s = r\rho_w/(\rho_m - \rho_w)$. Air, yang volumenya
tidak berubah, harus mengisi kedalaman yang sama di atas luas yang sama, sehingga
$r + s = \qty{574}{m}$, $r\,\rho_m/(\rho_m - \rho_w) = 574$, dan $r = 574\times2300/3300 =
\qty{400}{m}$: ratusan meter, hanya dari laju penyebaran.
\end{solution}

\begin{solution}{exo:b1:rifts-ridges-margins:11}
(a) $\sqrt{x - a} = \sqrt x\sqrt{1 - a/x} \approx \sqrt x\,(1 - a/2x)$ untuk
$a \ll x$, sehingga
$\Delta d \approx C\sqrt{x/u}\cdot a/(2x) = Ca/(2\sqrt{ux})$.
(b) Pada $x = \qty{450}{km}$: eksak \qty{1205}{m}, pendekatan
$350\times350/(2\sqrt{10.6\times450}) = \qty{887}{m}$; pada \qty{1350}{km}:
\qty{550}{m} dan \qty{512}{m}; pada \qty{3350}{km}: \qty{334}{m} dan
\qty{325}{m}. Rasio lompatan pendekatan terhadap eksak hanya bergantung pada
$r = a/x$: $(r/2)/(1 - \sqrt{1 - r})$, yang bernilai 0.9 untuk $r = 0.36$. Pendekatan
baik hingga \qty{10}{\percent} di atas $x = 350/0.36 \approx
\qty{970}{km}$, sekitar tiga kali pergeseran dari punggungan.
\end{solution}

\begin{solution}{exo:b1:rifts-ridges-margins:12}
(a) $d_2 - d_1 = 3671.1 - 3606.8 = \qty{64.3}{m}$ pada \qty{10}{Ma};
$4072.8 - 4065.3 = \qty{7.5}{m}$ pada 20; $4415.4 - 4417.0 = \qty{-1.7}{m}$ pada
30; $4707.5 - 4713.6 = \qty{-6.1}{m}$ pada 40. (b) $d_2 - d_1$ positif pada 20
dan negatif pada 30, dan kedua hukum kontinu: keduanya sama di antaranya.
Metode bagi dua: pada 25, $d_2 - d_1 = +0.9$; pada 27.5, $-0.7$; pada 26.25, $+0.01$; pada
26.875, $-0.3$; pada 26.56, $-0.2$; pada 26.41, $-0.08$; pada 26.33, $-0.04$: akar
terletak antara 26.25 dan 26.33, pada \qty{26.3}{Ma}, tempat keduanya memberi
\qty{4294}{m}. (c) $d_1' = 175/\sqrt{26.3} = \qty{34.1}{m.Ma^{-1}}$ dan
$d_2' = (3200/62.8)\,e^{-26.3/62.8} = \qty{33.5}{m.Ma^{-1}}$. Kedua hukum memiliki
nilai yang hampir sama dan \omterm{def:b1:geological-maps:attitude}{kemiringan} yang sama di sana, sehingga kurva gabungan tidak
memiliki tikungan yang terlihat.
\end{solution}

\begin{solution}{pb:b1:rifts-ridges-margins:1}
\textbf{1.} Pergeseran horizontal $= $ lemparan$/\tan$(\omterm{def:b1:geological-maps:attitude}{kemiringan}): F1 \qty{346}{m}, F2 \qty{770}{m},
F3 \qty{2014}{m}, F4 \qty{163}{m}, F5 \qty{462}{m}, F6 \qty{1172}{m}. F3, dengan
lemparan dan pergeseran horizontal terbesar, adalah sesar batas.

\textbf{2.} Perpanjangannya adalah jumlahnya, \qty{4927}{m}, sekitar \qty{4.9}{km}. Celah
itu selebar $60 - 4.9 = \qty{55.1}{km}$ sebelum pensesaran.

\textbf{3.} Normal terhadap celah yang berarah \ang{020} menunjuk ke \ang{110}; gerak ke
arah \ang{105} membentuk \ang{5} dengannya. Komponen tegak lurusnya adalah
$6.0\cos\ang{5} = \qty{5.98}{mm.yr^{-1}}$.

\textbf{4.} $\qty{4.927}{km}/\qty{5.98}{km.Ma^{-1}} = \qty{0.82}{Ma}$.

\textbf{5.} Dalam \qty{3.2}{Ma} lempeng-lempeng telah saling menjauh sejauh
$5.98\times3.2 = \qty{19.1}{km}$ melintasi celah. Sesar-sesar membawa
$4.93/19.1 = \qty{26}{\percent}$ darinya.

\textbf{6.} Tiga perempat lainnya, \qty{14.2}{km}, tidak berada pada sesar-sesar penampang
itu. Ia mungkin ditampung oleh dike yang menyusup sepanjang sumbu, yang membuka celah
tanpa mensesarkannya, seperti di Afar, atau oleh sesar di luar penampang. Dengan dike
selebar \qty{8}{m}, $\qty{14200}{m}/\qty{8}{m} \approx 1780$ dike diperlukan,
satu setiap $3.2\times10^6/1780 \approx 1800$ tahun rata-rata.

\textbf{7.} Dari C2An: $(21.0 + 20.2)/2.581 = 41.2/2.581 =
\qty{15.96}{mm.yr^{-1}}$. Dari C3n: $(41.0 + 40.2)/5.235 = 81.2/5.235 =
\qty{15.51}{mm.yr^{-1}}$.

\textbf{8.} $(41.2 - 2.2)/2.581 = \qty{15.11}{mm.yr^{-1}}$ dan
$(81.2 - 2.2)/5.235 = \qty{15.09}{mm.yr^{-1}}$: setelah dikoreksi, keduanya sepakat hingga
\qty{0.02}{mm.yr^{-1}}. Koreksi lebih berarti bagi kron yang lebih muda, yang jaraknya
lebih pendek.

\textbf{9.} Pergeseran keluar menggeser kedua tepi satu sisi dengan jumlah yang sama,
sehingga ia saling meniadakan dalam selisihnya. Untuk kedua sisi,
$(81.2 - 41.2)/(5.235 - 2.581) = 40.0/2.654 = \qty{15.07}{mm.yr^{-1}}$.

\textbf{10.} Timur laut: $(41.0 - 21.0)/2.654 = \qty{7.54}{mm.yr^{-1}}$;
barat daya: $(40.2 - 20.2)/2.654 = \qty{7.54}{mm.yr^{-1}}$. Penyebarannya simetris; selisih
\qty{0.8}{km} antara kedua sisi, sama untuk kedua kron, adalah pergeseran konstan, seperti
sumbu yang diletakkan \qty{0.4}{km} di luar pusat.

\textbf{11.} $\cos\Delta = \sin\ang{30.9}\sin\ang{18.0} +
\cos\ang{30.9}\cos\ang{18.0}\cos\ang{16.4} = 0.15870 + 0.78287 = 0.94157$, maka
$\Delta = \ang{19.69}$. Dengan $\omega R = 0.403\times(\pi/180)\times6371 =
\qty{44.81}{mm.yr^{-1}}$, lajunya adalah $44.81\sin\ang{19.69} =
\qty{15.10}{mm.yr^{-1}}$, laju terkoreksi pada soal 8 dan 9.

\textbf{12.} $\cos\Delta = \sin\ang{30.9}\sin\ang{26.0} +
\cos\ang{30.9}\cos\ang{26.0}\cos\ang{11.4} = 0.22512 + 0.75600 = 0.98112$, maka
$\Delta = \ang{11.15}$ dan lajunya adalah $44.81\sin\ang{11.15} =
\qty{8.7}{mm.yr^{-1}}$.

\textbf{13.} Prediksi, 15.1 dan \qty{8.7}{mm.yr^{-1}}, cocok dengan 15 dan
\qty{9}{mm.yr^{-1}} yang terukur. Laju itu adalah $\omega R\sin\Delta$, dan kutubnya
terletak di barat laut Laut Merah, di timur laut Libya: makin ke selatan, makin jauh dari
kutub, makin cepat pembukaannya.

\textbf{14.} $15.09\times5.235 = \qty{79}{km}$ dasar samudra.

\textbf{15.} $5000/120.6 = \qty{41}{km.Ma^{-1}}$, \qty{41}{mm.yr^{-1}} rata-rata.

\textbf{16.} $5000/15.1 = \qty{331}{Ma}$ pada laju tengah;
$5000/8.7 = \qty{575}{Ma}$ pada laju utara.

\textbf{17.} Membuka sekitar kutub yang berdekatan, Laut Merah melebar seperti sepasang
gunting: lebarnya bertambah sebanding dengan $\sin\Delta$, dan ia tetap berupa baji yang
menyempit ke arah kutub. Pada \qty{8.7}{mm.yr^{-1}}, bagian utara memerlukan
$79/8.7 = \qty{9.1}{Ma}$ untuk membuat \qty{79}{km} dasar samudra yang kini dimiliki
bagian tengah.

\textbf{18.} Pada \qty{15}{mm.yr^{-1}} sumbu Laut Merah adalah punggungan sangat lambat;
lembah aksial di antara sesar normal diharapkan.

\textbf{19.} $2500 + 350\times1 = \qty{2850}{m}$;
$2500 + 350\sqrt{10} = \qty{3607}{m}$; $2500 + 350\sqrt{40} = \qty{4714}{m}$.

\textbf{20.} $t = (1500/350)^2 = \qty{18.4}{Ma}$.

\textbf{21.} Dalam \qty{10}{Ma} pertama ia tenggelam $350\sqrt{10} = \qty{1107}{m}$;
antara 70 dan \qty{80}{Ma}, $350(\sqrt{80} - \sqrt{70}) = 350\times(8.944 -
8.367) = \qty{202}{m}$, sekitar lima kali lebih sedikit dalam waktu yang sama.

\textbf{22.} $\mathrm dd/\mathrm dt = C/(2\sqrt t) = 175/\sqrt{80} =
\qty{19.6}{m.Ma^{-1}}$, sekitar \qty{0.02}{mm} setahun.

\textbf{23.} $490/\sqrt{80} = \qty{54.8}{mW.m^{-2}}$.

\textbf{24.} Hukum akar kuadrat: $2500 + 350\sqrt{80} = 2500 + 350\times8.944 =
\qty{5630}{m}$. Hukum lempeng: $6400 - 3200\,e^{-80/62.8} = 6400 - 3200\times0.2798
= \qty{5505}{m}$. Keduanya berbeda sebesar \qty{126}{m}. Pada \qty{80}{Ma} kedua hukum
masih berlaku: hukum akar kuadrat hingga sekitar \qty{80}{Ma}, hukum lempeng dari sekitar
\qty{20}{Ma}; melampaui itu, hukum akar kuadrat menjadi terlalu dalam.

\textbf{25.} Hukum akar kuadrat memberi \qty{5630}{m}: hingga persepuluh kilometer
terdekat, dasar samudra baru terletak pada kedalaman
\textbf{\qty[text-series-to-math=true]{5.6}{km}} ketika berumur 80 juta tahun
(hukum lempeng memberi \qty{5.5}{km}).
\end{solution}
